% ============================================================================
%  SECCIÓN 5.1: ANÁLISIS DE LOS DATOS QUE NECESITAN PERSISTENCIA
% ============================================================================

\section{An\'alisis de los datos que necesitan persistencia}

En el avance anterior (PDF~4), la aplicaci\'on procesaba datos y respond\'ia a
las acciones del usuario, pero la informaci\'on cr\'itica se conservaba
\'unicamente en memoria: al cerrar la aplicaci\'on o reiniciar el dispositivo,
se perd\'ia la sesi\'on del usuario, los favoritos guardados y los escenarios
creados. En este avance se identifican los datos que requieren persistencia y
se implementan los mecanismos necesarios para conservarlos m\'as all\'a del
ciclo de vida de la aplicaci\'on.

\subsection{Identificaci\'on de datos cr\'iticos}

A partir del an\'alisis del c\'odigo fuente y de las funcionalidades
implementadas, se identificaron los siguientes conjuntos de datos que
requieren persistencia:

\begin{table}[H]
  \centering
  \caption{Datos que requieren persistencia}
  \contenidoConFuente{%
    \begin{tablaAPA}
      \begin{tabular}{@{}p{3.0cm}p{4.0cm}p{3.5cm}p{4.0cm}@{}}
        \toprule
        \textbf{Dato} & \textbf{Descripci\'on} & \textbf{Frecuencia de uso} &
        \textbf{Consecuencia de p\'erdida} \\
        \midrule
        Sesi\'on del usuario & Token JWT y datos del usuario autenticado &
        Continua & El usuario debe iniciar sesi\'on nuevamente. \\
        \addlinespace
        Favoritos del usuario & Lista de escenarios guardados por el usuario &
        Peri\'odica & El usuario pierde sus escenarios favoritos. \\
        \addlinespace
        Escenarios deportivos & Cat\'alogo de escenarios con sus detalles,
        im\'agenes y deportes & Lectura frecuente &
        La app no muestra datos sin conexi\'on. \\
        \addlinespace
        Eventos pr\'oximos & Actividades programadas en cada escenario &
        Consulta espor\'adica & No se muestran eventos futuros. \\
        \addlinespace
        Perfil del usuario & Nombre, email y rol del usuario &
        Consulta peri\'odica & No se muestra la informaci\'on del perfil. \\
        \bottomrule
      \end{tabular}
    \end{tablaAPA}%
  }{Elaboraci\'on propia en base al an\'alisis del c\'odigo fuente del
  proyecto.}
\end{table}

\subsection{Clasificaci\'on por criticidad}

Los datos identificados se clasifican en tres niveles de criticidad:

\begin{enumerate}
  \item \textbf{Cr\'iticos (alta criticidad):} datos cuya p\'erdida afecta
    directamente la experiencia del usuario. Incluyen la sesi\'on del usuario
    y los favoritos. Sin estos datos, la aplicaci\'on pierde funcionalidad
    esencial.

  \item \textbf{Importantes (media criticidad):} datos que mejoran la
    experiencia pero cuya ausencia es manejable. Incluyen el perfil del
    usuario y los eventos pr\'oximos. La aplicaci\'on puede funcionar sin
    ellos, aunque con informaci\'on incompleta.

  \item \textbf{De referencia (baja criticidad):} datos que se consultan
    frecuentemente pero que se obtienen del servidor. Incluyen el cat\'alogo
    de escenarios y sus relaciones. La aplicaci\'on puede mostrar datos
    parciales o un estado vac\'io mientras se recuperan del servidor.
\end{enumerate}

\subsection{An\'alisis de persistencia por componente}

A continuaci\'on se describe qu\'e datos persiste cada componente de la
aplicaci\'on y por qu\'e raz\'on:

\subsubsection{Servicio de Supabase (\textit{supabase.ts})}

El servicio de Supabase es el punto de entrada principal para todas las
operaciones de persistencia. Configura el cliente con las opciones de
persistencia de la sesi\'on:

\begin{lstlisting}
import AsyncStorage from
  '@react-native-async-storage/async-storage';
import { createClient } from '@supabase/supabase-js';

export const supabase = createClient(
  supabaseUrl, supabaseAnonKey, {
  auth: {
    storage: AsyncStorage,
    autoRefreshToken: true,
    persistSession: true,
  },
});
\end{lstlisting}

Las opciones configuradas son:

\begin{itemize}
  \item \texttt{storage: AsyncStorage} -- motor de persistencia local
    para los tokens de autenticaci\'on.
  \item \texttt{autoRefreshToken: true} -- renueva autom\'aticamente
    el token antes de que expire.
  \item \texttt{persistSession: true} -- guarda la sesi\'on en el
    almacenamiento local para sobrevivir cierres de la app.
\end{itemize}

\subsubsection{Contexto de autenticaci\'on (\textit{AuthContext.tsx})}

El contexto de autenticaci\'on gestiona la sesi\'on del usuario. Al
inicializarse, recupera la sesi\'on almacenada localmente y escucha
cambios de autenticaci\'on en tiempo real:

\begin{lstlisting}
useEffect(() => {
  supabase.auth.getSession()
    .then(({ data: { session } }) => {
      setSession(session);
      setLoading(false);
    });

  const { data: { subscription } } =
    supabase.auth.onAuthStateChange(
      (event, session) => {
        if (event === 'SIGNED_IN'
            || event === 'SIGNED_OUT') {
          queryClient.clear();
        }
        setSession(session);
      }
    );

  return () => subscription.unsubscribe();
}, []);
\end{lstlisting}

\subsubsection{Hooks de datos (\textit{useScenarios.ts},
\textit{useFavorites.ts}, \textit{useEvents.ts})}

Cada hook de datos consulta Supabase y almacena el resultado en el
cach\'e de React Query con un tiempo de validez de 5 minutos:

\begin{lstlisting}
return useQuery({
  queryKey: ['scenarios'],
  queryFn: fetchScenarios,
  staleTime: 1000 * 60 * 5, // 5 minutos
});
\end{lstlisting}

Este cach\'e en memoria permite que las consultas repetitivas no
realicen llamadas redundantes al servidor, mejorando el rendimiento
y reduciendo el consumo de datos m\'oviles.
